\chapter{Wstęp}
\label{chap:wstep}

Przedmiotem dokumentacji jest zaprojektowanie, wdrożenie oraz kompleksowa weryfikacja maszyny wirtualnej o nazwie \textbf{VulnBox}. Głównym założeniem projektu jest zbudowanie deterministycznego, izolowanego środowiska opartego na systemie operacyjnym \textbf{Ubuntu Server 26.04.1 LTS}, które demonstruje wieloetapowy łańcuch błędów konfiguracyjnych prowadzący od braku wiedzy początkowej o systemie (\texttt{0-knowledge}) do pełnych uprawnień administracyjnych (\texttt{root}).

\section{Cel i zakres projektu}
\label{sec:cel-zakres}

Projekt obejmuje warstwę inżynieryjną samej maszyny wirtualnej eksportowanej do formatu \texttt{VulnBox.ova} \cite{ovf20} oraz zestaw zautomatyzowanych skryptów wdrożeniowych i audytowych utrzymywanych w repozytorium projektu \cite{vulnbox_repo}. Architektura \texttt{VulnBox} została zaprojektowana tak, aby każdy etap weryfikacji wynikał bezpośrednio z błędów konfiguracji konkretnych usług systemowych lub aplikacyjnych, przy jednoczesnym zachowaniu pełnej odporności samego systemu bazowego na nieautoryzowane modyfikacje z poziomu konsoli wirtualizatora.

W ramach realizacji projektu zdefiniowano następujące zadania:
\begin{enumerate}
    \item Przygotowanie bazowego obrazu systemu \textbf{Ubuntu Server 26.04.1 LTS} (\texttt{amd64}) w hiperwizorze Oracle VirtualBox pod nazwą \texttt{VulnBox} (nazwa hosta \texttt{vulnbox}).
    \item Skonfigurowanie warstwy sieciowej z wykorzystaniem \textbf{Netplan} w sposób zapewniający bezobsługowe pobieranie adresacji IP przez protokół DHCP niezależnie od identyfikatora karty sieciowej po imporcie pliku \texttt{VulnBox.ova}.
    \item Zabezpieczenie \textbf{GRUB2} hasłem \texttt{PBKDF2-SHA512} przy zachowaniu automatycznego startu systemu (flaga \texttt{--unrestricted}) oraz wdrożenie ochrony z użyciem zapory \textbf{UFW} i usługi \textbf{Fail2Ban}.
    \item Wdrożenie czterowarstwowego łańcucha usług: zasobu plikowego \textbf{Samba} (\texttt{:445/tcp}), portalu diagnostycznego \textbf{Nginx + PHP-FPM} (\texttt{:8080/tcp}), wewnętrznego serwisu \textbf{Flask API} (\texttt{127.0.0.1:5000}) oraz mechanizmów administracyjnych (\texttt{cron} i \texttt{sudoers.d}).
    \item Przeprowadzenie pełnej weryfikacji manualnej (\emph{End-to-End}) oraz automatycznej przy użyciu skryptu audytowego \texttt{verify\_vm.sh}, połączonej z procedurą czyszczenia obrazu (\texttt{cleanup\_vm.sh}) przed eksportem do pliku \texttt{VulnBox.ova}.
\end{enumerate}

\section{Wymagania projektowe i założenia}
\label{sec:wymagania-konstrukcyjne}

W tabeli \ref{tab:wymagania-bazowe} zestawiono kluczowe wymagania stawiane maszynie wirtualnej \texttt{VulnBox} wraz z przyjętym w projekcie sposobem ich technicznej realizacji.

\begin{table}[H]
    \centering
    \caption{Zestawienie wymagań konstrukcyjnych dla maszyny wirtualnej VulnBox}
    \label{tab:wymagania-bazowe}
    \tablefont
    \begin{tabularx}{\textwidth}{L{3.8cm} L{4.2cm} Y}
        \toprule
        \hcell{Obszar wymagań} & \hcell{Parametr docelowy} & \hcell{Zastosowane rozwiązanie techniczne} \\
        \midrule
        System operacyjny & Ubuntu Server 26.04.1 LTS (bez GUI) & Instalacja minimalna bez pakietów X11/Wayland, zarządzanie przez \texttt{systemd} \\
        Wirtualizacja i eksport & Oracle VirtualBox, format \texttt{VulnBox.ova} & Standard OVF 2.0, pojedynczy adapter sieci izolowanej (\texttt{vboxnet0}) \\
        Konfiguracja sieci & Automatyczne DHCP (\emph{out-of-the-box}) & Reguła Netplan dopasowująca interfejsy \texttt{en*}/\texttt{eth*} z \texttt{dhcp-identifier: mac} \\
        Ochrona rozruchu & Blokada trybu \emph{Single User Mode} & Hasło \texttt{password\_pbkdf2} w \texttt{/etc/grub.d/40\_custom} i flaga \texttt{--unrestricted} \\
        Ochrona brzegowa & Filtracja portów i blokada siłowa SSH & Zapora UFW (\texttt{deny incoming}, otwarte TCP 22, 445, 8080) oraz profil \texttt{sshd} w Fail2Ban \\
        Ścieżka weryfikacji & Wielostopniowa (\texttt{0-knowledge} $\rightarrow$ \texttt{root}) & Łańcuch: Samba $\rightarrow$ Nginx/PHP $\rightarrow$ Flask Local API $\rightarrow$ konto \texttt{dev} $\rightarrow$ konto \texttt{root} \\
        Znaczniki etapów & Flagi weryfikacyjne \texttt{user.txt} i \texttt{root.txt} & Pliki z uprawnieniami \texttt{0600} w katalogach \texttt{/home/dev/user.txt} oraz \texttt{/root/root.txt} \\
        \bottomrule
    \end{tabularx}
\end{table}

\section{Struktura repozytorium projektowego}
\label{sec:struktura-repo}

Wszystkie kody źródłowe, konfiguracje usług oraz dokumentacja projektu są utrzymywane w publicznym repozytorium GitHub (\url{https://github.com/alehshehda/APNWA-VulnBox}) \cite{vulnbox_repo} podzielonym na dwa główne moduły oraz potoki CI/CD:
\begin{itemize}
    \item \texttt{report/} - źródła niniejszej dokumentacji projektowej w systemie \LaTeX{} kompilowane automatycznie przez potok \texttt{compile.yml} do pliku \texttt{APNWA-VulnBox.pdf} publikowanego w GitHub Releases oraz GitHub Pages,
    \item \texttt{box/} - kompletny pakiet wdrożeniowy maszyny wirtualnej \texttt{VulnBox}, obejmujący dokumentację instalacyjną \texttt{box/SETUP.md}, przewodnik weryfikacyjny \texttt{box/WALKTHROUGH.md}, pliki konfiguracyjne usług (\texttt{box/configs/}), kod źródłowy komponentów (\texttt{box/src/}) oraz skrypty automatyzujące budowę, weryfikację i czyszczenie systemu (\texttt{box/scripts/}), walidowane dodatkowo w potoku \texttt{box-ci.yml}.
\end{itemize}
